% !TeX root = ../final-report.tex
% TODO(team): Keep design details, policy formulas, and implementation boundaries aligned with the code.
\section{Design and Implementation}
\label{sec:design}
\input{assets/architecture}

\subsection{Control and Data Planes}
Figure~\ref{fig:architecture} separates metadata from segment delivery.
The Go backend comprises Controllers, storage Peers, a headless client, and a
Demo API gateway; a React/TypeScript frontend provides browser playback and
management. The evaluated deployment uses three Controllers and three Peers.
Controllers use HashiCorp Raft, BoltDB-backed persistent storage, and file
snapshots. A deterministic finite-state machine (FSM) applies metadata
commands; network I/O and clocks remain in handlers and the coordinator.
Reads verify leadership and execute a Raft barrier before exposing state.
Snapshots preserve manifests, registrations, tasks, and operation-ID
deduplication results. A write requires a majority; an uncertain response is
retried with the same operation ID.

\subsection{Verified Publication and Storage}
Ingestion creates a \texttt{DRAFT} manifest listing ordered segment IDs,
durations, sizes, and SHA-256 hashes, then uploads each segment directly to
two distinct Peers. A Peer checks size and hash, synchronizes a temporary
file, atomically installs its content-addressed blob, synchronizes the
directory, and durably replaces its index before returning a receipt.
Conflicting overwrites are rejected; equivalent retries recheck the stored
blob. Unindexed orphan bytes are retained and counted against capacity,
but are not advertised as replicas.

The Controller checks receipts against the configured Peer endpoints and
commits registrations through Raft. Let $\mathcal{S}_v$ be a video's segments
and $R_s$ its distinct registered holders. Publication commits
\texttt{DRAFT}$\rightarrow$\texttt{READY} only if
\begin{equation}
  \forall s\in\mathcal{S}_v,\qquad |R_s|\geq 2.
  \label{eq:publication}
\end{equation}
This is a check of confirmed registrations at publication, not a continuous
liveness guarantee. Files installed without a committed registration are
not returned as playback locations. Fixed placement is used in the study.
An optional resource-aware planner ranks eligible Peers by normalized free
space and upload capacity, with greater bandwidth weight for the first three
segments. Its capacity deductions are local to a plan; concurrent imports
have no global durable reservation, so the Peer quota is the final guard.

\subsection{Playback Paths}
The headless engine fetches published metadata and downloads segments with a
bounded prefetch window and worker pool. Round-robin rotates candidate
holders by segment index. Failed or invalid transfers try another holder;
exhausted candidates trigger a fresh location query and one further round.
Only size- and hash-verified segments become consumable. Consumption follows
manifest durations after a configurable contiguous startup buffer; payloads
are discarded rather than decoded. An optional deadline strategy ranks
sources using session-local throughput estimates, outstanding requests, and
failure penalties; the reported repair study holds round-robin fixed.

Browser playback instead fetches verified segments through the gateway and
checks them again with Web Crypto. FFmpeg encodes once into continuous
H.264/AAC fragmented MP4, preserving decode timestamps across segments.
The browser serializes append/remove operations on one MediaSource and
SourceBuffer, with bounded prefetch and historical buffers. Seeking cancels
obsolete requests and refills from the relevant keyframe. Browser waiting
events and headless simulated stalls are separate measurements. Metadata
caching and bounded retries bridge some Controller outages; uncached requests
still require an available quorum-backed service.

\subsection{Persistent Repair and Reconciliation}
Peers report heartbeat/inventory observations every two seconds. The leader
also polls configured endpoints independently; inventories rehash indexed
objects, excluding missing or corrupt data. Observations are ephemeral,
whereas replica registrations and tasks are replicated state. For segment
$s$, let $V_s(t)\subseteq R_s$ denote holders verified as reachable and valid
by the current inventory check. When $|V_s|<2$, the coordinator starts a
local grace timer. After grace, exactly one valid source and a distinct
eligible target are needed to create a committed \texttt{pending} task.
No source means no reconstructible copy. Grace restarts after leadership
loss and does not cancel tasks already created.

The coordinator resumes pending work before planning new work and processes
one segment per reconciliation cycle. It GETs source bytes, verifies them,
rechecks leadership, and PUTs the target with the task ID as operation ID.
A matching durable receipt allows one FSM command to register the target
and mark the task \texttt{completed} atomically. Copy errors can mark a task
\texttt{failed}; later cycles may plan another attempt. At most one task per
segment is pending. On resumption, an already-installed matching target
receipt avoids recopying, even if the source has stopped. Verified inventory
can also recover registration after a lost acknowledgment. An old leader
may install extra immutable bytes, but cannot publish them without quorum.

\subsection{Repair Bandwidth Controls}
All segment GETs share a Peer-wide token bucket (rate $C_p$, burst
32\,KiB). Repair reads additionally pass a bucket at rate
$\alpha_p(t)C_p$. The common token budget over $\Delta$ is
\begin{equation}
 \mathcal{B}_p(\Delta)=C_p\Delta+32\,\mathrm{KiB}.
 \label{eq:shared-budget}
\end{equation}
S sets $\alpha=1$;
F25 and F50 set $\alpha=0.25$ and $0.50$.
N and OFF are no-failure and failure-without-repair controls, respectively.
L updates every 250\,ms. Let $F_p$ count active foreground requests, $I_p$
be elapsed time since foreground activity became idle, and $A_p$ be the
oldest committed task age among active repair requests at that source
(zero if none):
\begin{equation}
 \begin{aligned}
 \beta_p(t)&=\begin{cases}
 1,&F_p(t)=0\ \text{and }I_p(t)\geq2\,\mathrm{s},\\
 0.25,&\text{otherwise},
 \end{cases}\\
 \alpha_p^{L}(t)&=\max\!\left\{\beta_p(t),\,
          0.5\,\mathbf{1}[A_p(t)>60\,\mathrm{s}]\right\}.
 \end{aligned}
 \label{eq:load-aware}
\end{equation}
Age begins at task creation, not at the fault; the 60\,s branch was not
activated in the main performance study. L observes source-local request
activity, not browser buffer levels. These allowances are additional caps,
not reserved shares or fairness guarantees, and token bursts and outstanding
chunk reservations affect short intervals. Counters record bytes written
at the HTTP response layer by traffic class, excluding protocol overhead
and TCP retransmissions.

The prototype also provides persistent video-deletion and surplus-replica
cleanup tasks and a local fault console. These support operation and testing;
the study concerns minimum-two repair in a trusted network, without
authentication or Byzantine protection.

\subsection{Implementation Validation}
The compact-study archive retains successful static checks, race-enabled
Go tests, and system checks for ingestion, repair, cleanup, deletion, and
media delivery. These correctness checks complement the E1--E5 trials.

\paragraph{Publication and object integrity.}
The metadata tests first attempt publication without enough replicas, then
confirm the same holder twice. Both publication attempts remain rejected;
only confirmation of a second distinct Peer permits publication. Repeating
the successful operation returns its original result, and snapshot restore
retains the published state. Storage tests interrupt an upload before its
declared size is received and check that no partial object enters inventory.
They also corrupt an installed blob and check that inventory excludes it
and an equivalent PUT cannot return a successful acknowledgment for the
corrupted bytes. Restart checks retain valid objects, while unindexed
residual blobs consume quota without becoming advertised copies. These
checks support distinguishing physical bytes, verified inventory, and
committed replica registrations in the recovery metric.

\paragraph{Repair interruption and acknowledgment recovery.}
The repair system check stops a holder, waits for a committed pending task,
and stops the Controller leader. It checks that the successor finishes the
same task, then stops the other original holder and verifies bytes from the
replacement by SHA-256. Further leader changes preserve the completed task.
A separate scenario installs a valid object on the replacement without
confirming it in metadata, modeling a lost upload response. With both
original holders stopped, inventory reconciliation discovers and registers
the installed copy; another leader failure retains that registration.
FSM tests reject competing pending tasks and wrong-target receipts, retain
pending state through snapshot restore, and accept repeated matching
completion. These checks exercise the gap between installing bytes and
committing a discoverable replica.

\paragraph{Bandwidth and media paths.}
HTTP-level tests transfer known payloads as repair, foreground, and probe
traffic, checking per-class byte counts and the unchanged common rate cap.
A separate foreground transfer checks that it cannot bypass that cap.
Policy tests cover the two-second idle transition and the strict
greater-than-60-second age boundary, which the main study does not activate.
Media acceptance checks ingest four real
H.264/AAC fragments, stop a Peer, verify gateway-delivered bytes, and decode
both the individual fragments and the joined stream. They also check
consistent initialization and continuous stored decode timelines. Earlier
Chromium acceptance separately exercised the MSE path and playback controls.
These checks establish tested delivery and decoding, without extending the
performance results to browser QoE or physical power loss.
